\section{Conclusion and Recommendation}
\label{sec:conclusion}

\subsection{Synthesis of Scientific Rationale}

The literature reviewed in this proposal collectively establishes that LBD is a molecularly complex, multifactorial disease whose individual biological layers---genetic risk variants \cite{chia2025gwas}, epigenetic dysregulation \cite{nabais2024methylation}, peripheral metabolic alterations \cite{frontiers2024metabolomics}, and systemic immune dysregulation \cite{singlecell2024immune}---each provide meaningful but incomplete insight into disease pathogenesis. The convergence of evidence across layers on specific biological themes---lysosomal biology centred on \textit{GBA1}, synaptic vesicle trafficking at the \textit{SYT16}/\textit{SNCA} loci, and neuroinflammation reflected in both methylation (UBASH3B) and immune cell activation---strongly supports the hypothesis that integration will yield mechanistic insight and biomarker power beyond single-layer analyses.

The existing comprehensive biomarker review \cite{boiten2025biomarkers} explicitly identifies integrated multi-marker panels as the most promising next step for the field, noting that no single protein, metabolite, or genetic marker is yet sufficient for accurate early LBD diagnosis. This proposal responds directly to that stated need.

\subsection{Assessment of Innovation and Impact}

The proposed project fills a genuine scientific gap. An integrative multi-omics analysis of LBD at the scale and methodological rigour proposed here has not yet been published. The expected outputs---a harmonised open-access dataset, a validated multi-omics biomarker panel, a convergent pathway map, and a prioritised drug target list---address practical needs in both the research and clinical communities. The fully open-source, containerised pipeline will have lasting value as infrastructure for the broader LBD research field.

\subsection{Potential Limitations to Acknowledge}

The most significant limitation is sample size. LBD remains an underrepresented condition in large-scale omics initiatives, and the statistical power for certain analyses (particularly those requiring matched multi-layer data within the same individuals) will be modest. Additionally, the majority of available data derives from populations of European ancestry, limiting the generalisability of findings. These limitations should be acknowledged transparently in any resulting publications and should be used to advocate for expanded, diversity-inclusive LBD multi-omics cohort development.

\subsection{Recommendation}

\textbf{This project is recommended for initiation.} The scientific rationale is strong, the methodology is sound and uses mature established tools, the required data are largely publicly accessible at no direct cost, and the computational infrastructure requirements are moderate and achievable. The project carries meaningful potential to generate high-impact findings---a clinically useful biomarker panel and biologically credible drug targets---within a reasonable timeframe of 12--18 months.

The recommended starting point is to focus initially on the three most data-rich omics layers (genomics via PRS, brain or blood DNA methylation, and bulk RNA-seq transcriptomics) using AMP-PD and GEO datasets, integrating them with MOFA+ to establish the pipeline end-to-end. Metabolomics and single-cell immunomics can be incorporated as additional layers in a second phase, contingent on successful data access requests and pilot feasibility. This phased approach reduces upfront risk while preserving the project's full scientific ambition.

Initiating a formal data access application to AMP-PD and preparing a targeted literature review for specific GEO dataset selection are the recommended immediate next steps.
